\chapter{الأشكال والفضاء}\label{ch:g1:shapes}

النوافذ مستطيلات، والعجلات دوائر، والسقوف تخفي مثلثات: الأشكال في
كل مكان. في هذا الفصل نتعلم كيف نعرفها بعدّ أضلاعها ورؤوسها، وكيف
نقول بدقة \emph{أين} تقع الأشياء: يمينًا، يسارًا، فوق، تحت،
بين.

\section{الأشكال الأربعة الأساسية}

\begin{definition}[المربع والمستطيل والمثلث والدائرة]\label{def:g1:shapes:def}
\begin{itemize}
  \item \emph{المثلث}\index{مثلث} له $3$ أضلاع و $3$ رؤوس؛
  \item \emph{المربع}\index{مربع} له $4$ أضلاع متساوية و $4$ رؤوس؛
  \item \emph{المستطيل}\index{مستطيل} له $4$ أضلاع و $4$ رؤوس،
        وضلعاه المتقابلان متساويان (مربع ممدود)؛
  \item \emph{الدائرة}\index{دائرة} مستديرة تمامًا: لا أضلاع لها ولا رؤوس.
\end{itemize}
\end{definition}

\begin{omfigure}
\begin{tikzpicture}[scale=0.7]
  \draw[thick, fill=omDef!12] (0,0) -- (2,0) -- (1,1.7) -- cycle;
  \node[below, align=center, font=\small] at (1,-0.25)
    {مثلث:\\$3$ أضلاع};
  \begin{scope}[xshift=3.6cm]
  \draw[thick, fill=omThm!12] (0,0) rectangle (1.7,1.7);
  \node[below, align=center, font=\small] at (0.85,-0.25)
    {مربع:\\$4$ أضلاع متساوية};
  \end{scope}
  \begin{scope}[xshift=6.9cm]
  \draw[thick, fill=omMeth!12] (0,0.05) rectangle (2.8,1.65);
  \node[below, align=center, font=\small] at (1.4,-0.25)
    {مستطيل:\\$4$ أضلاع};
  \end{scope}
  \begin{scope}[xshift=11.3cm]
  \draw[thick, fill=omProp!12] (0.9,0.85) circle (0.9);
  \node[below, align=center, font=\small] at (0.9,-0.25)
    {دائرة:\\بلا رؤوس};
  \end{scope}
\end{tikzpicture}

{\small الأشكال الأربعة الأساسية. ولتسمية شكل، لا تثق بالنظرة
الأولى: \emph{عُدّ} أضلاعه ورؤوسه.}
\end{omfigure}

\begin{example}[الشكل المائل يبقى كما هو]\label{ex:g1:shapes:tilted}
المربع الواقف على أحد رؤوسه يبقى مربعًا --- أربعة أضلاع متساوية
وأربعة رؤوس --- ولو بدا كأنه طائرة ورقية. وإدارة الشكل لا تغيّر ما
هو.
\end{example}

\begin{example}[صيد الأشكال]\label{ex:g1:shapes:hunt}
في الصف: السبورة مستطيل، والساعة دائرة، والمنديل المطويّ قد يكون
مثلثًا، وبعض النوافذ مربعات. والأشياء الحقيقية ليست أشكالًا كاملة،
لكن وجوهها قريبة منها (وتعود المجسمات ووجوهها في
\cref{ch:g2:shapes}).
\end{example}

\section{أين هو؟}

\begin{definition}[كلمات الموضع]\label{def:g1:shapes:position}
لنقول أين يقع الشيء: \emph{يمين} / \emph{يسار}،
\emph{فوق} / \emph{تحت}، \emph{على} / \emph{أسفل}،
\emph{أمام} / \emph{خلف}، \emph{بين}،
\emph{داخل} / \emph{خارج}.
\end{definition}

\begin{omfigure}
\begin{tikzpicture}[scale=0.75]
  \draw[thick, fill=omDef!12] (0,0) circle (0.55);
  \draw[thick, fill=omThm!12] (1.9,-0.55) rectangle (3.0,0.55);
  \draw[thick, fill=omMeth!12] (4.4,-0.5) -- (5.5,-0.5) -- (4.95,0.45)
    -- cycle;
  \node[below, font=\small] at (0,-0.75) {دائرة};
  \node[below, font=\small] at (2.45,-0.75) {مربع};
  \node[below, font=\small] at (4.95,-0.75) {مثلث};
  \node[right, font=\small, align=left] at (6.3,0)
    {المربع بين\\الدائرة والمثلث؛\\
    والدائرة يسار المربع};
\end{tikzpicture}

{\small كلمات الموضع في العمل. واِنتبه: يسارك أنت ويسار من
يقابلك متعاكسان!}
\end{omfigure}

\section{المسارات على الشبكة}

\begin{method}[وصف مسار]\label{met:g1:shapes:path}
على الورق المربّعات، المسار قائمة خطوات:
$\to$ (مربع واحد يمينًا)، $\leftarrow$ (يسارًا)، $\uparrow$ (فوق)،
$\downarrow$ (تحت).
\begin{enumerate}
  \item لاتّباع مسار: اِبدأ من المربع المعلَّم وأطِع الأسهم واحدًا
        واحدًا؛
  \item ولوصف مسار: اِمشِ فيه بذهنك واكتب الأسهم بالترتيب؛
  \item ويمكن لمسارين مختلفين أن يصلا بين المربعين نفسيهما ---
        أحدهما أقصر والآخر أطول.
\end{enumerate}
\end{method}

\begin{omfigure}
\begin{tikzpicture}[scale=0.6]
  \draw[gray!40] (0,0) grid (7,4);
  \fill[omDef!40] (0.5,0.5) +(-0.42,-0.42) rectangle +(0.42,0.42);
  \node[font=\footnotesize] at (0.5,0.5) {البداية};
  \fill[omThm!40] (5.5,2.5) +(-0.42,-0.42) rectangle +(0.42,0.42);
  \node[font=\footnotesize] at (5.5,2.5) {الهدف};
  % seven unit arrows, one per step, with small gaps
  \foreach \a/\b in {0.5/1.5, 1.5/2.5}
    \draw[-Stealth, omDef, very thick] (\a+0.1,0.5) -- (\b-0.1,0.5);
  \foreach \a/\b in {0.5/1.5, 1.5/2.5}
    \draw[-Stealth, omDef, very thick] (2.5,\a+0.1) -- (2.5,\b-0.1);
  \foreach \a/\b in {2.5/3.5, 3.5/4.5}
    \draw[-Stealth, omDef, very thick] (\a+0.1,2.5) -- (\b-0.1,2.5);
  \draw[-Stealth, omDef, very thick] (4.6,2.5) -- (5.05,2.5);
\end{tikzpicture}

{\small مسار من البداية إلى الهدف، سهم لكل خطوة:
$\to\ \to\ \uparrow\ \uparrow\ \to\ \to\ \to$ --- سبع خطوات. هل
تجد مسارًا آخر بالعدد نفسه من الخطوات؟}
\end{omfigure}

\section{تمارين}

\begin{exercise}[$\star$]\label{exo:g1:shapes:1}
لكل شكل، عُدّ أضلاعه ورؤوسه ثم سمِّه: شكل له $3$ رؤوس؛ وشكل
مستدير بلا رؤوس؛ وشكل له $4$ أضلاع
متساوية.
\end{exercise}

\begin{exercise}[$\star$]\label{exo:g1:shapes:2}
اِبحث في البيت عن: مستطيلين، ودائرة، ومثلث، ومربع.
(الأبواب؟ الأطباق؟ لافتات الطريق؟)
\end{exercise}

\begin{exercise}[$\star$]\label{exo:g1:shapes:3}
صواب أم خطأ؟ «المربع المُدار على أحد رؤوسه يكفّ عن كونه
مربعًا». اِشرح مستعينًا بما في \cref{ex:g1:shapes:tilted}.
\end{exercise}

\begin{exercise}[$\star$]\label{exo:g1:shapes:4}
اُرسم على الورق المربّعات: مربعًا طوله $3$ مربعات في $3$؛ ومستطيلًا
$5$ في $2$؛ ومثلثًا يتخذ قطر الشبكة ضلعًا له.
\end{exercise}

\begin{exercise}[$\star$]\label{exo:g1:shapes:5}
ضع ثلاث لعب في صف. وصِف الصف بكلمات الموضع: أيّها بين الأخريين؟
وأيّها على اليسار؟
\end{exercise}

\begin{exercise}[$\star$]\label{exo:g1:shapes:6}
اُرسم دائرة. ثم اُرسم علامة × \emph{داخلها}، ونجمة
\emph{خارجها}، ونقطة \emph{على} الدائرة نفسها.
\end{exercise}

\begin{exercise}[$\star$]\label{exo:g1:shapes:7}
على الورق المربّعات، علّم مربع البداية. واتبع:
$\to\ \uparrow\ \to\ \uparrow\ \to$. وعلّم موضع الوصول. ثم اُكتب
المسار الذي يعود مباشرة إلى البداية.
\end{exercise}

\begin{exercise}[$\star$]\label{exo:g1:shapes:8}
كم رأسًا في المجموع: مثلثان ومربع؟ ودائرة
ومستطيلان؟
\end{exercise}

\begin{exercise}[$\star\star$]\label{exo:g1:shapes:9}
خذ $12$ عود ثقاب. هل تستطيع أن ترسم بها \emph{كلها} حدود مربع،
بأعواد كاملة وبالعدد نفسه على كل ضلع؟ كم عودًا في كل ضلع؟ وهل يمكن
ذلك مع $10$ أعواد؟
\end{exercise}
